\documentclass[10pt, aspectratio=169]{beamer}

\usetheme{default}
\usecolortheme{default}

\definecolor{petazul}{RGB}{29, 82, 122}    
\definecolor{petlaranja}{RGB}{241, 90, 36}

\setbeamercolor{structure}{fg=petazul}
\setbeamercolor{title}{fg=petazul}
\setbeamercolor{frametitle}{fg=petazul}

% Margem laranja
\setbeamertemplate{headline}{
  \begin{beamercolorbox}[wd=\paperwidth,ht=3pt,dp=0pt]{}
    \color{petlaranja}\rule{\paperwidth}{3pt}
  \end{beamercolorbox}
}

\usepackage[utf8]{inputenc}
\usepackage[portuguese]{babel}
\usepackage{graphicx}
\usepackage{listings}
\usepackage{xcolor}
\usepackage{tikz}

\lstset{
    backgroundcolor=\color{gray!5},   
    basicstyle=\ttfamily\tiny\color{black},
    breaklines=true,
    keywordstyle=\color{petazul}\bfseries,
    commentstyle=\color{gray}\itshape,
    stringstyle=\color{black},
    showstringspaces=false,
    extendedchars=true,
    literate={á}{{\'a}}1 {ã}{{\~a}}1 {é}{{\'e}}1 {í}{{\'i}}1 {ó}{{\'o}}1 {õ}{{\~o}}1 {ç}{{\c{c}}}1 {Ú}{{\'U}}1
}

%[Infos da apa]
\title{Desenvolvimento de Webpages Dinâmicas}
\subtitle{Catálogo de Atividades da Engenharia UFF}
\author{Gabriel Maia Ramos}
\institute{
    Universidade Federal Fluminense - UFF \\
    Escola de Engenharia (TCE) \\
    Departamento de Engenharia de Telecomunicações (TET)
}
\date{\today}

\begin{document}

% 1[Capa]
\begin{frame}
    \titlepage
\end{frame}

% [Sumário]
\begin{frame}{Agenda da Apresentação}
    \tableofcontents
\end{frame}


\section{Introdução e Motivação}


\begin{frame}{Contextualização}
    \textbf{Cenário}
    \begin{itemize}
        \item A UFF conta com dezenas de grupos (Empresas Juniores, Coletivos, Projetos).
        \item As informações estavam dispersas, incompletas, ou desatualizadas em páginas antigas.
    \end{itemize}
    
    \vspace{0.5cm}
    \textbf{Proposta de Solução}
    \begin{itemize}
        \item Desenvolver uma plataforma web unificada e dinâmica, capaz de centralizar essas informações de forma padronizada.
        \item Garantir fácil manutenção e escalabilidade dos dados sem a necessidade de reescrever o código.
    \end{itemize}
\end{frame}

\begin{frame}{Motivações do Projeto}
    \begin{itemize}
        \item \textbf{Busca e Tratamento de Dados:} O alto volume de grupos gerou um processo exaustivo de coleta, filtragem e correção manual das informações.
        \vspace{0.3cm}
        \item \textbf{Experiência do Usuário:} Transformar dados brutos e fragmentados em um catálogo interativo, responsivo e de navegação fluida.
        \vspace{0.3cm}
        \item \textbf{Aperfeiçoamento Técnico:} Aplicar o Modelo Cliente-Servidor na prática, estudando a comunicação entre JSON e PHP/HTML.
    \end{itemize}
\end{frame}


\section{Fundamentação e Arquitetura}


\begin{frame}{Fundamentação Teórica: Modelo Cliente-Servidor}
    \centering
    \begin{tikzpicture}
        \node[draw, thick, minimum width=3.5cm, minimum height=1cm] (C) at (0,0) {Cliente (Navegador)};
        \node[draw, thick, minimum width=3.5cm, minimum height=1cm] (S) at (7,0) {Servidor (Apache/PHP)};
        
        \draw[->, thick] (C.north east) -- (S.north west) node[midway, above, font=\scriptsize] {Requisição HTTP};
        \draw[<-, thick] (C.south east) -- (S.south west) node[midway, below, font=\scriptsize] {Resposta HTML Dinâmica};
    \end{tikzpicture}
    
    \vspace{1cm}
    \begin{itemize}
        \item \textbf{Front-end (HTML/CSS):} Estrutura semântica e \textit{design} da interface.
        \item \textbf{Back-end (PHP):} Motor lógico que lê e processa os dados no servidor.
        \item \textbf{Base de Dados (JSON):} Estrutura leve de intercâmbio de dados.
    \end{itemize}
\end{frame}

\begin{frame}{Arquitetura e Modularização}
    O projeto foi estruturado utilizando o princípio de \textbf{modularização}, dividindo as responsabilidades para manter o código limpo:
    
    \vspace{0.5cm}
    \begin{itemize}
        \item \textbf{\texttt{dados/}}: Diretório isolado contendo os dados brutos (\texttt{.json}).
        \item \textbf{\texttt{style.css}}: Regras visuais e "animações".
        \item \textbf{\texttt{index.php}}: Trabalha \textbf{unido ao HTML}, processando a lógica e inserindo dados dinânmicos na estrutura da página.
    \end{itemize}
\end{frame}


\section{Lógica de Processamento}


\begin{frame}[fragile]{Lógica 1: Memória}
\begin{lstlisting}[language=php]
// Mapeamento modular das categorias
$categorias = [
    "Empresas Juniores" => "dados/empresas_juniores.json",
    // ...
];

// Validacao defensiva e Alocacao em Memoria
if (file_exists($caminho_arquivo)) {
    $json_puro = file_get_contents($caminho_arquivo);
    $grupos = json_decode($json_puro, true);
}
\end{lstlisting}
    
    \vspace{0.2cm}
    \textbf{Como funciona:} Um vetor organiza quais arquivos devem ser lidos. A função \texttt{file\_exists()} previne falhas, e \texttt{json\_decode()} traduz o texto para a memória do PHP,
\end{frame}

\begin{frame}[fragile]{Lógica 2: Tratamento de Inconsistências}
\begin{lstlisting}[language=php]
foreach ($grupos as $grupo) {
    // Tratamento de espacos invisiveis e inconsistencias manuais
    $site = isset($grupo['website']) ? trim($grupo['website']) : '';

    if (!empty($site) && $site != "Nao informado") {
        // Inserindo o protocolo de rede caso o usuario tenha esquecido
        $link = (strpos($site, 'http') !== 0) ? "https://" . $site : $site;
        $classe = "card"; // Elemento ganha interatividade no CSS
    } else {
        $link = "#";
        $classe = "card sem-link"; // Elemento tem sua interatividade desativada
    }
    // O PHP injeta essas variaveis ($link, $classe) direto no HTML gerado
}
\end{lstlisting}
    \vspace{0.2cm}
    \textbf{Como funciona:} Esse bloco converte erros humanos (URLs incompletas, termos como "Não possui") em uma interface visualmente correta e livre de \textit{links} quebrados.
\end{frame}


\section{Interface e Conclusão}


\begin{frame}{Decisões de Design (UI/UX)}
    \begin{itemize}
        \item \textbf{Identidade Visual:} Adoção da paleta de cores institucional do PET-Tele.
        \item \textbf{Carrossel Horizontal:} Inspirado em vitrines digitais, permite a exibição de múltiplos cartões (\textit{Cards}), o que aproveita a altura da tela.
        \item \textbf{Interações:} Efeitos de transição (\textit{hover}) elevam os cartões ativos, enquanto cartões sem \textit{site} cadastrado mantêm-se estáticos, oferecendo uma página limpa e clara ao usuário.
    \end{itemize}
\end{frame}

\begin{frame}{Conclusões}
    \begin{itemize}
        \item \textbf{Objetivo Alcançado:} Transição de um cenário estático para um ecossistema dinâmico.
        \item \textbf{Manutenibilidade:} A atualização do catálogo pode ser feita por qualquer membro do grupo através de edições em texto (JSON), sem risco de quebra da interface.
        \item \textbf{Escalabilidade:} O código modular permite a fácil adição de novas categorias.
    \end{itemize}
\end{frame}

\begin{frame}
    \centering
    \Large \textbf{Obrigado!} \\
    \vspace{1cm}
    \normalsize Perguntas? \\
    \vspace{1cm}
    \small gabrielmaia@id.uff.br
\end{frame}

\end{document}